\subsection{空间与单连通空间的积}

命题 8. —— 设 B 为拓扑空间，T 为单连通且局部连通的空间，E 为 B × T 的覆叠，其投影为 p，并设 t 为 T 的一点。以 \(E_{t}\) 表示空间 \(p^{-1}(B \times \{t\})\) ，它带有映射 \(p_{t} = pr_{1} \circ p \mid E_{t}: E_{t} \to B\) ；这是 B 的一个覆叠。则存在唯一的从覆叠 ( \(E_{t} \times T, p_{t} \times Id_{T}\) ) 到覆叠 E 上的 (B × T)-同构，它把每个 \(x \in E_{t}\) 的 (x, t) 映为 x 。

可以假设 B 非空。设 \(x\) 为 \(E_t\) 的一点。把 I, 第 125 页命题 3 应用于覆叠 E 和连续映射 T → B × T, \(u \mapsto (p_t(x), u)\) ，便存在唯一的连续映射 \(f_x \colon \mathrm{T} \to \mathrm{E}\) ，使得 \(f_x(t) = x\) 且 \(p(f_x(u)) = (p_t(x), u)\) 对所有 \(u \in T\) 成立。设 \(h \colon E_t \times T \to E\) 为由 \(h(x, u) = f_x(u)\) 定义的映射。我们有 \(h(x, t) = x\) 与 \(p \circ h = p_t \times \operatorname{Id}_T\) 。映射 h 是映射 \(p_t \times \operatorname{Id}_T\) 到 E 中的提升。h 在 \(E_t \times \{t\}\) 上的限制是连续的，h 在 \(\{x\} \times T\) 上的限制对 \(E_t\) 的每一点 x 也是连续的。由于空间 T 局部连通，映射 h 连续（I, 第 37 页, 定理 1 的推论 1）。

设 \(b\) 为 B 的一点。按构造，映射 \(h\) 诱导从纤维 \(p_t^{-1}(b) \times \{t\}\) （\(E_t \times T\) 的纤维）到纤维 \(p^{-1}(b,t)\) （E 在点 \((b,t)\) 处的纤维）上的双射。由于空间 T 连通且局部连通，映射 \(h\) 是双射（I, 第 84 页, 命题 7 的推论）。因此它是 \(B \times T\) -同构（I, 第 30 页, 命题 6 的推论 2）。

设 \(h'\) 为一个 \((\mathrm{B} \times \mathrm{T})\) -同构，它从覆叠 \((\mathrm{E}_{t} \times \mathrm{T}, p_{t} \times \mathrm{Id}_{\mathrm{T}})\) 映到覆叠 E 上，且把 \((x, t)\) 映为 x（对每一点 \(x \in E_{t}\) ）。对所有 \(x \in E_{t}\) ，映射 \(u \mapsto h(x, u)\) 与 \(u \mapsto h'(x, u)\) 相等（I, 第 34 页, 命题 11 的推论 1）。于是有 \(h = h'\) 。

推论 1. —— 在命题 8 的假设下，若 t 与 t' 为 T 的两点，则覆叠 \(E_{t}\) 与 \(E_{t'}\) 是 B-同构的。

推论 2. —— 在命题 8 的假设下，设 (E,p) 与 (E',p') 为 B × T 的覆叠，并设 \(k: E_t \to E_t'\) 为 B-态射。则存在唯一的 (B × T)-态射 \(\widetilde{k}\) : E → E'，它扩张 k。若 k 是 B-同构，则 \(\widetilde{k}\) 是 (B × T)-同构。

由命题 8，可以假设存在 B 的覆叠 F 与 F′，使得 E 与 E′ 分别为 (B × T)-空间 F × T 与 F′ × T。于是有 \(E_{t} = F \times \{t\}\) ，\(E_{t}' = F' \times \{t\}\) ，且映射 k 写作 \((x, t) \mapsto (k'(x), t)\) ，其中 \(k'\) 是从 F 到 F' 中的 B-态射。映射 \(k' \times Id_{T}\) 是扩张 k 的 B × T-态射；若 k 是同构，则它也是同构。

设 \(\widetilde{k}\colon\mathrm{E}\to\mathrm{E}^{\prime}\) 为 \((\mathrm{B}\times\mathrm{T})\) -态射，它扩张 \(k\) 。设 \(x\) 为 F 的一点。以 \(q\) 表示 F 的投影，并令 \(b=q(x)\) 。映射 \(u\mapsto\widetilde{k}(x,u)\) 与 \(u\mapsto(k^{\prime}(x),u)\) 是在 \(\mathrm{E}^{\prime}\) 中的提升，所提升的映射为 \(u\mapsto(b,u)\) ，它把 T 映到 \(\mathrm{B}\times\mathrm{T}\) 。它们在 \(t\) 处相合。由于空间 T 连通，它们相等。于是，\(\widetilde{k}\) 就是 \((\mathrm{B}\times\mathrm{T})\) -态射 \(k^{\prime}\times\mathrm{Id}_{\mathrm{T}}\colon\mathrm{F}\times\mathrm{T}\to\mathrm{F}^{\prime}\times\mathrm{T}\) 。

推论 3. —— 设 B 与 B' 为拓扑空间，T 为单连通且局部连通的拓扑空间。设 \(f: B' \times T \to B\) 为连续映射，E 为 B 的覆叠。给定 T 的一点 t 和一个连续提升 \(g_t: B' \times \{t\} \to E\) （\(f | B' \times \{t\}\) 到 E 中的提升），存在唯一的 f 到 E 中的连续提升 g，它扩张 \(g_t\) 。

由 I, 第 9 页命题 3，只需证明 \(f^{*}(\mathrm{E})\) 在 \(\mathrm{B}' \times \{t\}\) 上的每个连续截面都唯一地扩张为 \(f^{*}(\mathrm{E})\) 的连续截面。而这由把推论 2 应用于覆叠 \((\mathrm{B}' \times \mathrm{T}, \mathrm{Id}_{\mathrm{B}' \times \mathrm{T}})\) 与 \(f^{*}\mathrm{E}\) （二者都是 \(\mathrm{B}' \times \mathrm{T}\) 的覆叠）即得。

注. —— 保留命题 8 的假设与记号。设 G 为离散拓扑群。假设 E 是以 G 为群的主覆叠。若给覆叠 \(E_{t}\) （B 的）与覆叠 \(E_{t} \times T\) （\(B \times T\) 的）赋予由 E 的结构导出的以 G 为群的主覆叠结构（I, 第 92 页, 例 1 与例 4），则覆叠 E 与 \(E_{t} \times T\) 是同构的主覆叠。事实上，设 \(h: E_{t} \times T \to E\) 为唯一的 (B × T)-态射，使得 \(h(x, t) = x\) （对所有 \(x \in E_{t}\) ）。对 G 的每个元素 g，映射 \((x, u) \mapsto h(x \cdot g, u) \cdot g^{-1}\) 是把 \((x, t)\) （对所有 \(x \in E_{t}\) ）映为 x 的 (B × T)-态射，因此等于 h。这就证明了 h 是主覆叠的 (B × T)-态射。

特别地，在前述假设之下，主覆叠 \(E_{t}\) 与 \(E_{t'}\) （对 \(t' \in T\) ）是同构的。类似地可以证明：若在推论 2 中 E 与 \(E'\) 是以 G 为群的主覆叠，且 k 是以 G 为群的主覆叠的 B-同构，则 \(\widetilde{k}\) 是主覆叠的 \((\mathrm{B} \times \mathrm{T})\) 同构。

命题 9. —— 设 B 为拓扑空间，(E,p) 为 B 的覆叠。设 T 为单连通、局部连通且局部紧的拓扑空间。以 \(\widetilde{p}\colon\mathcal{C}_{c}(\mathrm{T};\mathrm{E})\to\mathcal{C}_{c}(\mathrm{T};\mathrm{B})\) 表示映射 \(g\mapsto p\circ g\) 。设 t 为 T 的一点。以 \(e_{\mathrm{E}}\colon\mathcal{C}_{c}(\mathrm{T};\mathrm{E})\to\mathrm{E}\) 表示把每个 \(g\in\mathcal{C}_{c}(\mathrm{T};\mathrm{E})\) 对应到 \(g(t)\) 的映射，以 \(e_{\mathrm{B}}\colon\mathcal{C}_{c}(\mathrm{T};\mathrm{B})\to\mathrm{B}\) 表示把每个 \(f\in\mathcal{C}_{c}(\mathrm{T};\mathrm{B})\) 对应到 \(f(t)\) 的映射。

方块

\[
\begin{array}{c} \mathcal {C} _ {c} (\mathrm{T}; \mathrm{E}) \xrightarrow {e _ {\mathrm{E}}} \mathrm{E} \\ \Big \downarrow \widetilde {p} \\ \mathcal {C} _ {c} (\mathrm{T}; \mathrm{B}) \xrightarrow {e _ {\mathrm{B}}} \mathrm{B} \end{array}
\]

是笛卡尔的。

我们先证明一个引理。

引理. —— 设 X, Y, Z 为拓扑空间，\(g: Y \rightarrow Z\) 为连续映射。

\begin{enumerate}
\def\labelenumi{\alph{enumi})}
\item
  映射 \(f\mapsto g\circ f\) （从 \(\mathcal{C}_c(\mathrm{X};\mathrm{Y})\) 到 \(\mathcal{C}_c(\mathrm{X};\mathrm{Z})\) ）是连续的。
\item
  若拓扑空间 Z 是豪斯多夫的，则映射 \(h \mapsto h \circ g\) （从 \(\mathcal{C}_{c}(\mathrm{Z};\mathrm{X})\) 到 \(\mathcal{C}_{c}(\mathrm{Y};\mathrm{X})\) ）是连续的。
\end{enumerate}

给定 X 的紧子集 K、Z 的开子集 U 以及从 X 到 Y 的连续映射 f，要有 \((g \circ f)(\mathrm{K}) \subset \mathrm{U}\) ，必须且只需 \(f(\mathrm{K}) \subset g^{-1}(\mathrm{U})\) 。因此第一个断言由紧收敛拓扑的定义得出（TG, X, 第 26 页, 定义 1）。

类似地，设 K 为 Y 的紧子集，U 为 X 的开子集。由于假设空间 Z 是豪斯多夫的，集合 \(g(\mathrm{K})\) 是紧的（TG, I, 第 63 页, 推论 1）。若 h 是从 Z 到 X 中的一个映射，则条件 \((h \circ g)(\mathrm{K}) \subset \mathrm{U}\) 无非就是条件 \(h(g(\mathrm{K})) \subset \mathrm{U}\) ，由此得出第二个断言。

现在来证明命题 9。由引理，映射 \(\widetilde{p}\) 连续；由 TG, X, 第 27 页的注 1，映射 \(e_{\mathrm{E}}\) 与 \(e_{\mathrm{B}}\) 连续。对任意映射 \(g \in \mathcal{C}_c(\mathrm{T};\mathrm{E})\) ，我们有

\[
(p \circ e _ {\mathrm{E}}) (g) = p (g (t)) = (p \circ g) (t) = \widetilde {p} (g) (t) = \left(e _ {\mathrm{B}} \circ \widetilde {p}\right) (g),
\]

于是命题中的方块图是交换的。

设 \(\varphi\colon\mathcal{C}_{c}(\mathrm{T};\mathrm{E})\to\mathcal{C}_{c}(\mathrm{T};\mathrm{B})\times_{\mathrm{B}}\mathrm{E}\) 为由 \(\varphi(g)=(p\circ g,g(t))\) （对所有 \(g\in\mathcal{C}_{c}(\mathrm{T};\mathrm{E})\) ）定义的连续映射。由 I, 第 125 页命题 3，它是双射。事实上，对每一对 \((f,x)\in\mathcal{C}_{c}(\mathrm{T};\mathrm{B})\times_{\mathrm{B}}\mathrm{E},\varphi^{-1}(f,x)\) 是到 E 中的唯一连续提升 \(g\) （它提升 \(f\) ），使得 \(g(t)=x\) 。由于空间 T 局部紧，映射 \(\psi\colon(\mathcal{C}_{c}(\mathrm{T};\mathrm{B})\times_{\mathrm{B}}\mathrm{E})\times\mathrm{T}\to\mathrm{B}\) （由 \(\psi((f,x),u)=f(u)\) 定义）是连续的（TG, X, 第 28 页, 推论 1）。由前面的推论 3，映射 \(\psi\) 有唯一的连续提升 \(\theta\colon(\mathcal{C}_{c}(\mathrm{T};\mathrm{B})\times_{\mathrm{B}}\mathrm{E})\times\mathrm{T}\to\mathrm{E}\) ，使得 \(\theta((f,x),t)=x\) （对 \((f,x)\in\mathcal{C}_{c}(\mathrm{T};\mathrm{B})\times_{\mathrm{B}}\mathrm{E}\) ）。于是有 \(\theta((f,x),u)=\varphi^{-1}(f,x)(u)\) （对 \((f,x)\in\mathcal{C}_{c}(\mathrm{T};\mathrm{B})\times_{\mathrm{B}}\mathrm{E}\) 与 \(u\in\mathrm{T}\) ）。由 TG, X, 第 28 页定理 3，映射 \((f,x)\mapsto\theta((f,x),\cdot)\) （从 \(\mathcal{C}_{c}(\mathrm{T};\mathrm{B})\times_{\mathrm{B}}\mathrm{E}\) 到 \(\mathcal{C}_{c}(\mathrm{T};\mathrm{E})\) ）是连续的，也就是说 \(\varphi^{-1}\) 连续。

这样，映射 \(\varphi\) 就是从 \(\mathcal{C}_c(\mathrm{T};\mathrm{E})\) 到纤维积 \(\mathcal{C}_c(\mathrm{T};\mathrm{B})\times_{\mathrm{B}}\mathrm{E}\) 上的同胚，由此即得命题（I, 第 8 页, 命题 2）。

